\documentclass[10pt,a4paper]{amsart}
\usepackage{amsmath,amssymb,amsthm,graphicx}
% Retained minimal preamble for the Stacks Project chapter more-algebra.tex (More on Algebra).
% The chapter file pinned-original-source.tex carries only \input{preamble} on its first line, so no
% preamble text lives inside the chapter itself.  The Stacks Project uses one single shared preamble for
% every chapter; the snapshot stages it identically as sources/stacks-chow/preamble.tex and
% sources/stacks-etale-cohomology/preamble.tex.  The block below is copied mechanically, line for line,
% from sources/stacks-chow/preamble.tex lines 14-18 (the Xy-pic setup) and 154-217 (the native theorem
% declarations together with the mathematical macros).  Nothing was added, removed or redefined; the
% wrapper class and the amsmath/amssymb/amsthm/graphicx base packages are supplied by the assembler.
% Chapter: more-algebra.tex -- More on Algebra.
% For commutative diagrams we use Xy-pic
\usepackage[all]{xy}

% We use 2cell for 2-commutative diagrams.
\xyoption{2cell}

% Theorem environments.
%
\theoremstyle{plain}
\newtheorem{theorem}[subsection]{Theorem}
\newtheorem{proposition}[subsection]{Proposition}
\newtheorem{lemma}[subsection]{Lemma}

\theoremstyle{definition}
\newtheorem{definition}[subsection]{Definition}
\newtheorem{example}[subsection]{Example}
\newtheorem{exercise}[subsection]{Exercise}
\newtheorem{situation}[subsection]{Situation}

\theoremstyle{remark}
\newtheorem{remark}[subsection]{Remark}
\newtheorem{remarks}[subsection]{Remarks}

\numberwithin{equation}{subsection}

% Macros
%
\def\lim{\mathop{\mathrm{lim}}\nolimits}
\def\colim{\mathop{\mathrm{colim}}\nolimits}
\def\Spec{\mathop{\mathrm{Spec}}}
\def\Hom{\mathop{\mathrm{Hom}}\nolimits}
\def\Ext{\mathop{\mathrm{Ext}}\nolimits}
\def\SheafHom{\mathop{\mathcal{H}\!\mathit{om}}\nolimits}
\def\SheafExt{\mathop{\mathcal{E}\!\mathit{xt}}\nolimits}
\def\Sch{\mathit{Sch}}
\def\Mor{\mathop{\mathrm{Mor}}\nolimits}
\def\Ob{\mathop{\mathrm{Ob}}\nolimits}
\def\Sh{\mathop{\mathit{Sh}}\nolimits}
\def\NL{\mathop{N\!L}\nolimits}
\def\CH{\mathop{\mathrm{CH}}\nolimits}
\def\proetale{{pro\text{-}\acute{e}tale}}
\def\etale{{\acute{e}tale}}
\def\QCoh{\mathit{QCoh}}
\def\Ker{\mathop{\mathrm{Ker}}}
\def\Im{\mathop{\mathrm{Im}}}
\def\Coker{\mathop{\mathrm{Coker}}}
\def\Coim{\mathop{\mathrm{Coim}}}

% Boxtimes
%
\DeclareMathSymbol{\boxtimes}{\mathbin}{AMSa}{"02}

%
% Macros for moduli stacks/spaces
%
\def\QCohstack{\mathcal{QC}\!\mathit{oh}}
\def\Cohstack{\mathcal{C}\!\mathit{oh}}
\def\Spacesstack{\mathcal{S}\!\mathit{paces}}
\def\Quotfunctor{\mathrm{Quot}}
\def\Hilbfunctor{\mathrm{Hilb}}
\def\Curvesstack{\mathcal{C}\!\mathit{urves}}
\def\Polarizedstack{\mathcal{P}\!\mathit{olarized}}
\def\Complexesstack{\mathcal{C}\!\mathit{omplexes}}
% \Pic is the operator that assigns to X its picard group, usage \Pic(X)
% \Picardstack_{X/B} denotes the Picard stack of X over B
% \Picardfunctor_{X/B} denotes the Picard functor of X over B
\def\Pic{\mathop{\mathrm{Pic}}\nolimits}
\def\Picardstack{\mathcal{P}\!\mathit{ic}}
\def\Picardfunctor{\mathrm{Pic}}
\def\Deformationcategory{\mathcal{D}\!\mathit{ef}}
\usepackage{cleveref}
\crefname{originallemma}{original lemma}{original lemmas}
\Crefname{originallemma}{Original lemma}{Original lemmas}
\makeatletter
\newlabel{algebra-lemma-cover}{{00EO}{source}{}{sourcepointer-1}{}}
\newlabel{algebra-lemma-cover@cref}{{[originallemma][0][]00EO}{[1][0][]source}{}{}{}}
\makeatother
\crefname{originallemma}{original lemma}{original lemmas}
\Crefname{originallemma}{Original lemma}{Original lemmas}
\makeatletter
\newlabel{lemma-invertible}{{36118}{source}{}{sourcepointer-2}{}}
\newlabel{lemma-invertible@cref}{{[originallemma][0][]36118}{[1][0][]source}{}{}{}}
\makeatother
\title{The Picard group of a unique factorization domain is trivial}
\author{The Stacks Project Authors}
\date{stacks-project:more-algebra.tex:local-snapshot-20261012; GNU FDL 1.2 or later}
\begin{document}
\maketitle
\setcounter{subsection}{1}
\noindent\textit{Editorial note.} All delivered mathematical text is the author's own, copied from the pinned Stacks Project chapter file more-algebra.tex; no statement or proof step was rewritten, reordered or removed.\par\medskip
\noindent\textit{Editorial note.} The chapter file pinned-original-source.tex carries only \textbackslash{}input\{preamble\} on its first line, so no preamble text lives inside the chapter. The retained file sources/stacks-more-algebra/retained-preamble.tex is a mechanical, line-for-line copy of the project's single shared preamble as staged in the snapshot; nothing was added, removed or redefined and no package was added.\par\medskip
\noindent\textit{Editorial note.} The source URL is the source-level Stacks Project URL, as required when more than one card is produced from one source; the per-claim tag is carried by claim\_key.\par\medskip
\noindent\textit{Editorial note.} The author writes 'g = v prod\_\{i in J\} p\_j\textasciicircum{}\{b\_j\}' using j both as the summation index and as the symbol introduced by the set J; the expression is reproduced unchanged.\par\medskip
\noindent\textit{Allowed standard prerequisites (editorial).}\begin{itemize}
\item Invertible modules are finite locally free of rank one and their duals are inverse. invoked at proof line 36189. Justifies that L is finite locally free and that varphi(L) is an invertible ideal, which is what makes the localization argument available.
\item Uniqueness of prime factorization in a UFD and its preservation under localization. used at proof lines 36201-36218. Is the arithmetic input that forces b\_j = c\_j in the final comparison.
\item The cover criterion for an ideal to be generated by a given element. invoked at proof line 36210. Reduces generating I by h to the case where I\_a is principal for a single a.
\end{itemize}\medskip
\ifdefined\hypertarget\hypertarget{sourcepointer-1}{}\fi
\noindent\textit{Source reference.} Original lemma 00EO: focused-seeds/AGP/sources/stacks-snapshot-20261011/tags.txt, source lines 542--542. Cross-chapter navigation pointer into the Stacks Project; the standard fact it supplies is applied inside the counted proof span and the argument resting on it stays there.\par\medskip
\ifdefined\hypertarget\hypertarget{sourcepointer-2}{}\fi
\noindent\textit{Source reference.} Original lemma 36118: sources/stacks-more-algebra/pinned-original-source.tex, source lines 36118--36118. Same-chapter navigation pointer; the statement it names is standard for the section and the argument that consumes it stays inside the counted proof span.\par\medskip
\begin{lemma}
\label{lemma-UFD-Pic-trivial}
Let $R$ be a UFD. Then $\Pic(R)$ is trivial.
\end{lemma}

\begin{proof}
Let $L$ be an invertible $R$-module. By Lemma \ref{lemma-invertible}
we see that $L$ is a finite locally free $R$-module. In particular
$L$ is torsion free and finite over $R$. Pick a nonzero element
$\varphi \in \Hom_R(L, R)$ of the dual invertible module.
Then $I = \varphi(L) \subset R$ is an ideal which is an invertible module.
Pick a nonzero $f \in I$ and let
$$
f = u p_1^{e_1} \ldots p_r^{e_r}
$$
be the factorization into prime elements with $p_i$ pairwise distinct.
Since $L$ is finite locally free there exist $a_i \in R$,
$a_i \not \in (p_i)$ such that $I_{a_i} = (g_i)$ for some $g_i \in R_{a_i}$.
Then $p_i$ is still a prime element of the UFD $R_{a_i}$ and
we can write $g_i = p_i^{c_i} g'_i$ for some $g'_i \in R_{a_i}$
not divisible by $p_i$. Since $f \in I_{a_i}$ we see that $e_i \geq c_i$.
We claim that $I$ is generated by $h = p_1^{c_1} \ldots p_r^{c_r}$ which
finishes the proof.

\medskip\noindent
To prove the claim it suffices to show that $I_a$ is generated by $h$
for any $a \in R$ such that $I_a$ is a principal ideal
(Algebra, Lemma \ref{algebra-lemma-cover}). Say $I_a = (g)$.
Let $J \subset \{1, \ldots, r\}$ be the set of $i$ such that
$p_i$ is a nonunit (and hence a prime element) in $R_a$. Because
$f \in I_a = (g)$ we find the prime factorization
$g = v \prod_{i \in J} p_j^{b_j}$
with $v$ a unit and $b_j \leq e_j$. For each $j \in J$ we have
$I_{aa_j} = g R_{aa_j} = g_j R_{aa_j}$, in other words
$g$ and $g_j$ map to associates in $R_{aa_j}$. By uniqueness
of factorization this implies that $b_j = c_j$ and the proof is complete.
\end{proof}

\end{document}
